\documentclass[11pt]{article}

\usepackage[T1]{fontenc}
\usepackage{lmodern}
\usepackage{amsmath,amssymb,amsthm}
\usepackage[margin=1in]{geometry}
\usepackage{array,booktabs,longtable}
\usepackage{enumitem}
\usepackage{microtype}
\usepackage[hidelinks]{hyperref}

\hypersetup{
  pdftitle={Mathematical errata and statement audit for three chemotaxis--logistic papers}
}

\newcommand{\R}{\mathbb{R}}
\newcommand{\Om}{\Omega}
\newcommand{\Tmax}{T_{\max}}
\newcommand{\norm}[1]{\left\lVert #1\right\rVert}
\newcommand{\abs}[1]{\left\lvert #1\right\rvert}
\newcommand{\False}{\textsc{False statement}}
\newcommand{\Gap}{\textsc{Proof gap}}
\newcommand{\Clarify}{\textsc{Clarification}}
\newcolumntype{L}[1]{>{\raggedright\arraybackslash}p{#1}}
\setlength{\emergencystretch}{3em}

\newtheorem{proposition}{Proposition}
\newtheorem{lemma}{Lemma}
\newtheorem{remark}{Remark}

\title{Mathematical Errata and Statement Audit\\
\large Three papers on chemotaxis systems with logistic sources}
\author{}
\date{27 July 2026}

\begin{document}
\maketitle

\begin{abstract}
This note gives a self-contained mathematical audit of three papers on
parabolic--elliptic chemotaxis systems.  Each finding is classified as an
explicitly false statement, a gap in a printed proof, or a clarification of
the intended continuation or regularity framework.  False statements are
accompanied by explicit counterexamples.  For proof gaps, the invalid step,
the missing hypothesis, and the range actually supported by the calculation
are stated separately.  Every conclusion below follows from the displayed
mathematical arguments.
\end{abstract}

\tableofcontents

\section{Sources and logical classification}

The audited sources are:
\begin{enumerate}[leftmargin=2em]
\item Wenxian Shen, \emph{Existence, uniqueness, stability, and monotonicity
  of traveling waves for repulsion/attraction chemotaxis models with logistic
  type source}, arXiv:2605.04401v1 (the ``traveling-wave paper'');
\item Le Chen, Ian Ruau, and Wenxian Shen,
  \emph{Chemotaxis models with signal-dependent sensitivity and a
  logistic-type source, I: Boundedness and global existence},
  arXiv:2512.14858v1 (``Part~I'');
\item Le Chen, Ian Ruau, and Wenxian Shen,
  \emph{Chemotaxis models with signal-dependent sensitivity and a
  logistic-type source, II: Persistence and stabilization},
  arXiv:2604.02599v1 (``Part~II'').
\end{enumerate}

\begin{description}[leftmargin=2.1em,style=nextline]
\item[\False] The printed hypotheses admit an explicit solution or test
  object for which the printed conclusion fails.
\item[\Gap] A displayed implication is invalid or a hypothesis required by
  an invoked estimate is missing.  This classification does not assert that
  the theorem itself is false.
\item[\Clarify] The mathematical intent is recoverable, but the literal
  statement is stronger than the continuation or regularity framework
  supports.
\end{description}

\begin{center}
\renewcommand{\arraystretch}{1.18}
\begin{longtable}{@{}L{0.05\textwidth}L{0.15\textwidth}L{0.44\textwidth}L{0.25\textwidth}@{}}
\toprule
\# & Source & Finding & Status\\
\midrule
\endhead
E1 & Traveling wave & Theorems 1.2--1.3: the Section 5 weighted-energy
  calculation contains a root obstruction, a frame mismatch, and four
  coefficient defects & \Gap\\
E2 & Traveling wave & Lemma 4.2: the asserted constant subsolution fails for
  an admissible member of the printed trap set & \False\\
E3 & Part I & Theorem 1.2 permits \(a>0,b=0\), for which spatially constant
  solutions grow exponentially & \False\\
E4 & Part I & Theorem 1.3(iv) invokes an elliptic estimate outside its
  exponent range & \Gap\\
E5 & Part I & Proposition 1.1 attaches an endpoint alternative to a finite
  horizon instead of the maximal continuation & \Clarify\\
E6 & Part I & Lemma 2.7 assumes an inequality only for \(t>0\) but concludes
  a bound including \(t=0\) & \False\ as written\\
E7 & Part II & Theorem 2.1(1) permits the pure-decay regime \(a=0<b\) &
  \False\\
E8 & Part II & Theorems 2.2--2.5 claim an all-time \(C^1\) estimate from
  continuous or merely \(L^\infty\)-small initial data & \Clarify\\
\bottomrule
\end{longtable}
\end{center}

\section{Traveling-wave paper}

\subsection{E1: the weighted stability calculation}

Let
\[
  \kappa=\frac{c-\sqrt{c^2-4}}2,
  \qquad \kappa^2-c\kappa+1=0.
\]
The printed stability interval begins at \(\eta=\kappa\).  The coefficient
produced by the global estimate in Section~5 is instead
\begin{equation}
  q(\eta)=\eta^2-(c-A)\eta+(1+B),                       \label{eq:q}
\end{equation}
where the displayed quantities in the paper have \(A\ge0\), and \(B>0\)
when \(\chi\ne0\).

\begin{proposition}[root obstruction]
Suppose \(A\ge0\), \(B>0\), and the discriminant of \(q\) is positive.  Then
\(\kappa\) is not in the negative set of \(q\).  If \(\kappa\) lies to the
left of the smaller root
\[
  \kappa^-=\frac{(c-A)-\sqrt{(c-A)^2-4(1+B)}}2,
\]
then \(q(\eta)>0\) for every \(\eta\) in a nonempty right-neighborhood of
\(\kappa\).
\end{proposition}

\begin{proof}
Using the defining equation for \(\kappa\),
\[
 q(\kappa)
 =(\kappa^2-c\kappa+1)+A\kappa+B
 =A\kappa+B>0.
\]
Since \(q\) is continuous, it remains positive near \(\kappa\).  When
\(\kappa<\kappa^-\), the quadratic is positive on
\((-\infty,\kappa^-)\), negative only between its two roots, and hence the
claimed neighborhood is contained in the positive region.
\end{proof}

Consequently, the displayed global coefficient estimate can prove decay only
for weights above \(\kappa^-\), not for every \(\eta>\kappa\).  This is a
gap in the proof: it does not rule out a different localized or spectral
argument recovering the larger interval.

\paragraph{Frame mismatch.}
After the change of variables \(z=x-ct\), the estimated energy is
\[
 E_{\rm mov}(t)=
 \int_\R e^{2\eta z}\abs{u(t,z+ct)-U(z)}^2\,dz.
\]
The quantity printed in (1.21) is
\[
 E_{\rm lab}(t)=
 \int_\R e^{2\eta x}\abs{u(t,x)-U(x-ct)}^2\,dx.
\]
The substitution \(x=z+ct\) gives the exact relation
\begin{equation}
 E_{\rm lab}(t)=e^{2\eta ct}E_{\rm mov}(t).             \label{eq:frames}
\end{equation}
Therefore decay of \(E_{\rm mov}\) does not by itself imply decay of the
printed laboratory-weighted quantity.  Either the conclusion must use the
co-moving weight \(e^{2\eta(x-ct)}\), or the decay rate must be strong enough
to absorb the factor in \eqref{eq:frames}.  There is also a sign error:
if \(\lambda=q(\eta)<0\), then \(e^{\lambda t}\), not \(e^{-\lambda t}\),
is the decaying exponential.

\paragraph{Coefficient audit.}
Writing \(w=u-U\) and \(z=v-V\), the elliptic identity
\(v_{xx}=v-u^\gamma\) gives the zero-order flux difference
\[
  (v a_m-a_{m+\gamma})w+U^m z.
\]
After multiplication by \(-\chi\), the perturbation equation contains
\[
 -\chi(v a_m-a_{m+\gamma}+b_2)w-\chi b_4z.
\]
Thus the signs of both \(a_{m+\gamma}\) and \(b_4z\) in (5.18)--(5.19)
must be reversed.  Absolute-value estimates can repair those two signs, but
three subsequent numerical coefficients must also be changed:
\begin{align}
 \abs{\chi b_1w_xw}
 &\le \frac12w_x^2+\frac12\abs{\chi}^2b_1^2w^2,
                                                               \label{eq:young}\\
 \int_\R V^2
 &\le\frac{\gamma^2M^{2(\gamma-1)}}{(1-\eta)^2}
       \int_\R W^2,                                           \label{eq:V}\\
 \int_\R V_x^2
 &\le\frac{\gamma^2M^{2(\gamma-1)}}{1-\eta^2}
       \int_\R W^2.                                           \label{eq:Vx}
\end{align}
Equation \eqref{eq:young} is Young's inequality with
\(a=w_x\) and \(b=\chi b_1w\).  The factor
\(M^{2(\gamma-1)}\) in \eqref{eq:V}--\eqref{eq:Vx} follows by squaring the
Lipschitz bound
\(\abs{s^\gamma-t^\gamma}\le\gamma M^{\gamma-1}\abs{s-t}\).
Finally, an upper bound on \(U_x/U\) does not imply an absolute-value bound;
the use in (5.23) requires a two-sided estimate for the logarithmic
derivative.

\paragraph{Range supported by the corrected calculation.}
After correcting the flux signs, \eqref{eq:young}--\eqref{eq:Vx}, and the
logarithmic-derivative estimate, the same global quadratic argument supports
the window
\[
  \kappa^-<\eta<
  \frac{1}{1+\abs{\chi}^{1/6}},
\]
provided this interval is nonempty and the corrected discriminant is
positive.  The energy must be written in the co-moving coordinate.  Any
claim for the larger interval beginning at \(\kappa\) needs an additional
argument not present in Section~5.

\subsection{E2: the constant subsolution in Lemma 4.2}

The lemma asserts that a sufficiently small constant \(W=d\) is a
subsolution for every profile in a printed trap set.  Consider
\[
 u(y)=\min\{1,e^{-y/2}\}\min\{1,\abs{y+1}\}.
\]
This function is continuous, nonnegative, satisfies the printed upper
barrier for \(\kappa=\tfrac12\), \(M=1\), and has
\[
 u(-1)=0,\qquad u(0)=1.
\]
For \(m=\alpha=\gamma=1\), its frozen elliptic response at \(-1\) is
\[
 V_0=\frac12\int_\R e^{-\abs{-1-y}}u(y)\,dy>0.
\]
Choose
\[
 \chi=\frac2{V_0},\qquad
 \kappa=\frac12,\quad \widetilde\kappa=1,\quad M=1,\quad c=\frac52,
\]
take \(D\) above the printed threshold, and choose a positive \(d\) below
the printed constant threshold.  At \(x=-1\), all derivatives of \(W=d\)
vanish, whereas the frozen operator is
\[
\begin{aligned}
 \mathcal A(d;u)(-1)
 &= -\chi d\bigl(V_0-u(-1)\bigr)+d(1-d)\\
 &= -2d+d-d^2=-d(1+d)<0.
\end{aligned}
\]
A subsolution requires \(\mathcal A(d;u)\ge0\).  The universal statement is
therefore false.

\begin{proposition}[a sufficient repaired condition]
Assume the trap satisfies \(0\le u\le M\), so its elliptic response obeys
\(0\le V\le M^\gamma\).  A constant \(d>0\) is a subsolution whenever
\begin{equation}
  \abs{\chi}\,d^{m-1}M^\gamma\le 1-d^\alpha.             \label{eq:repair}
\end{equation}
\end{proposition}

\begin{proof}
For a constant test function the diffusion and first-derivative terms
vanish.  The chemotactic zero-order contribution is bounded below by
\(-\abs{\chi}d^mM^\gamma\), while the logistic contribution is
\(d(1-d^\alpha)\).  Hence
\[
 \mathcal A(d;u)\ge
 d\bigl(1-d^\alpha-\abs{\chi}d^{m-1}M^\gamma\bigr),
\]
which is nonnegative under \eqref{eq:repair}.
\end{proof}

For example, when \(m=1\), the simple assumptions
\(\abs{\chi}M^\gamma\le\tfrac12\) and \(0<d^\alpha\le\tfrac12\)
imply \eqref{eq:repair}.

\section{Part I: boundedness and global existence}

\subsection{E3: Theorem 1.2 needs a reaction guard}

The printed assumptions \(a,b\ge0\) include \(a>0,b=0\).  On any smooth
bounded Neumann domain, take \(c_0>0\) and set
\[
 u(t,x)=c_0e^{at},\qquad
 v(t,x)=\frac{\nu}{\mu}u(t,x)^\gamma.
\]
All spatial derivatives and boundary fluxes vanish.  The elliptic equation
holds, and the parabolic equation reduces to \(u_t=au\).  Thus this is a
positive global classical solution for every \(m,\beta,\chi_0\), but
\[
 \norm{u(t)}_\infty=c_0e^{at}\longrightarrow\infty.
\]
The counterexample even satisfies the small-sensitivity condition by taking
\(\chi_0=0\).  Under the standing nonnegativity assumptions, the natural
repair is
\[
  a=0\quad\text{or}\quad b>0,
\]
which excludes exactly the branch \(a>0,b=0\).  The more conservative
case split used by the printed proof is
\((a=b=0)\lor(a>0\land b>0)\).

\subsection{E4: Theorem 1.3(iv) uses an inadmissible exponent}

Alternative (iv) assumes
\[
 \beta\ge\frac12,\qquad \alpha=2m+\gamma-2.
\]
Section~5.4 applies Proposition~2.2 with
\[
 s(P)=\frac{P+\alpha}{\gamma},
\]
although that proposition requires \(s(P)>1\).  Under the critical identity,
\[
 s(P)>1
 \quad\Longleftrightarrow\quad
 P>2-2m.
\]
The iteration begins at
\[
 q_*=\max\{1,N\alpha/2\};
\]
hence the printed route requires the missing condition
\begin{equation}
 \boxed{\max\{1,N\alpha/2\}>2-2m.}                     \label{eq:missing-exp}
\end{equation}

The printed hypotheses do not imply \eqref{eq:missing-exp}.  Indeed,
\[
 N=1,\quad m=\frac14,\quad \gamma=\frac72,\quad
 \alpha=2,\quad\beta=1
\]
satisfies \(\alpha=2m+\gamma-2\), but
\[
 q_*=1,\qquad s(q_*)=\frac67<1.
\]
Moreover \((N\alpha-2)_+=0\), so the first smallness disjunct in (1.25) is
automatic.  Thus Proposition~2.2 is invoked outside its stated domain.
Adding \eqref{eq:missing-exp} repairs this proof route; without it, a
different argument is required.

\subsection{E5: Proposition 1.1 must use the maximal horizon}

A finite endpoint alternative cannot hold for every finite local horizon.
When \(a,b>0\), the positive equilibrium
\[
 u(t,x)=\left(\frac ab\right)^{1/\alpha}
\]
is a classical solution on every finite interval and neither blows up nor
approaches zero at its endpoint.

The standard corrected statement is as follows.  Let \(\mathcal T\) be the
set of times to which the initial solution can be continued, and put
\[
 T_{\max}=\sup\mathcal T.
\]
Uniqueness makes the local solutions compatible on overlaps, so they glue to
a solution on \([0,T_{\max})\).  If \(T_{\max}<\infty\), the continuation
criterion implies that at least one controlling norm becomes unbounded or
the positive lower bound degenerates as \(t\uparrow T_{\max}\).  If neither
happened, local existence at a terminal slice would extend the solution
beyond \(T_{\max}\), contradicting its definition.  The endpoint alternative
therefore belongs to the distinguished maximal continuation, not to an
arbitrary finite restriction of it.

\subsection{E6: Lemma 2.7 omits the initial trace}

The lemma assumes a differential inequality only on \((0,\Tmax)\) but
concludes a supremum bound on \([0,\Tmax)\).  Let \(\Tmax=1\) and take
\[
 u(t,x)=t^{-1},\qquad 0<t<1.
\]
For \(p=2\), the mass quantity is \(Y(t)=t^{-2}\), and
\[
 Y'(t)=-2t^{-3}\le-t^{-3}.
\]
Thus the interior inequality is compatible with the lemma's nonnegative
constants, while
\[
 \sup_{0<t<1}Y(t)=\infty.
\]
No inequality imposed only at positive times can control an omitted initial
trace.

The repair is to assume \(Y\) extends continuously to \(t=0\), or at least
that \(Y(0)<\infty\) is part of the data.  Scalar comparison may then be
applied on \([0,T]\) for every \(T<\Tmax\), with the comparison constant
depending on \(Y(0)\) and the coefficients in the differential inequality.
The classical PDE initial trace supplies exactly this datum.

\section{Part II: persistence and stabilization}

\subsection{E7: Theorem 2.1(1) omits the pure-decay regime}

The printed assumptions allow \(a=0<b\).  For \(c_0>0\), define
\[
 u(t,x)=\bigl(c_0^{-\alpha}+\alpha bt\bigr)^{-1/\alpha},
 \qquad
 v(t,x)=\frac{\nu}{\mu}u(t,x)^\gamma.
\]
All spatial terms vanish and
\[
 u_t=-bu^{1+\alpha}.
\]
The solution is global, bounded, and strictly positive for every finite
time, yet
\[
 \lim_{t\to\infty}\inf_{x\in\Om}u(t,x)=0.
\]
The positive persistence conclusion is therefore false in this regime.  The
case split actually used in the proof is
\[
 (a=b=0)\quad\text{or}\quad(a>0\land b>0).
\]
When \(a>0,b=0\), a globally bounded positive solution cannot exist because
its mass satisfies
\[
 \frac{d}{dt}\int_\Om u=a\int_\Om u,
\]
so that branch of an implication restricted to bounded global solutions is
vacuous.

\subsection{E8: an all-time \(C^1\) estimate cannot follow from \(C^0\) data}

Equation (2.12) asserts a \(C^1\) exponential estimate for every \(t\ge0\),
whereas Theorem~2.2 assumes only continuous, positive initial data close in
\(L^\infty\).  On \(\Om=(0,1)\), consider the Neumann data
\[
 u_{0,N}(x)=u_*+N^{-1/2}\cos(N\pi x).
\]
For large \(N\) they are positive, have the same prescribed mass as \(u_*\),
and satisfy
\[
 \norm{u_{0,N}-u_*}_{L^\infty}=N^{-1/2}\to0,
 \qquad
 \norm{\partial_xu_{0,N}}_{L^\infty}=\pi N^{1/2}\to\infty.
\]
Thus no \(L^\infty\)-ball has a uniform \(C^1\) bound at \(t=0\).  For a
merely continuous datum, the \(C^1\) norm may not even be defined.  This is
consistent with parabolic smoothing, whose first-derivative estimate has the
short-time form
\[
 \norm{\partial_x e^{t\Delta}f}_\infty
 \le Ct^{-1/2}\norm f_\infty.
\]

There are two correct formulations:
\begin{enumerate}[leftmargin=2em]
\item if the initial perturbation is small in a fractional domain embedded
  into \(C^1\), sectorial stability gives an all-time strong-norm estimate;
\item for continuous or \(L^\infty\)-small data, first evolve to any suitable
  \(t_0>0\), then obtain exponential \(C^1\) convergence for \(t\ge t_0\).
\end{enumerate}
The linear stability threshold is unaffected; only the initial-time
regularity assertion changes.

\section{Conclusion}

Four findings, E2, E3, E6, and E7, are false as literally stated and are
settled by explicit counterexamples.  E1 and E4 identify precise gaps in
printed estimates without asserting that the associated final theorems are
false.  E5 and E8 are repaired by attaching the conclusions to the correct
maximal-continuation and positive-time regularity frameworks.  Keeping these
logical categories separate is essential: a defective proof step, an
overstated norm at \(t=0\), and a counterexample to the conclusion are
different mathematical claims.

\end{document}
